\section{Results}\label{sec:results}

Figure~\ref{fig:responses} and Table~\ref{tab:results} report, for each outcome and each policy, three objects:
the worst-case identified set of Proposition~\ref{prop:bounds} (grey), the effect on the overlap region
$\hat\theta^{\cA}_d$ with its 90\% bootstrap interval (blue), and the breakdown value $\lambda^\ast_d$ of
Proposition~\ref{prop:lambda}. Responses are cumulative changes from the month of the meeting: percent for
industrial production and the PCE price index, percentage points for the unemployment rate. The overlap region
covers 67\% of meetings for a rate increase (menus $\{U,T\}$ and $\{U,T,E\}$) and 59\% for a decrease ($\{U,E\}$ and
$\{U,T,E\}$). The main specification uses the Greenbook propensity score.

\begin{figure}[!htbp]
  \centering
  \includegraphics[width=0.92\linewidth]{responses_main.pdf}
  \caption{Cumulative effects of a rate increase (left) and decrease (right) relative to no change. Grey: worst-case
  identified set (Proposition~\ref{prop:bounds}). Blue: effect on the overlap region $\hat\theta^{\cA}_d$ with 90\%
  moving-block bootstrap interval. Greenbook propensity score; 133 meetings, 1989--2006.}\label{fig:responses}
\end{figure}

\paragraph{Worst-case bounds are not informative.} With $K$ equal to the range of outcomes the economy actually
experienced, every identified set contains zero at every horizon, and so does every \citet{imbens2004} interval.
Two years after a rate increase, the identified set for industrial production runs from $-7.2\%$ to $+4.3\%$. The
menus that exclude the action or the status quo account for a third or more of all meetings, and without an
assumption about them the data cannot sign the average effect. The 2020 version reported much narrower bounds
because it calibrated $K$ differently (Appendix~\ref{app:revision}).

\paragraph{The overlap-region effects have the wrong sign for real activity.} Where both the action and the status
quo were on the menu, the effects are point-identified and estimated with some precision. After a rate
\emph{increase}, industrial production is higher than under no change (by 1.5\% at six months and 1.9\% at twelve,
with 90\% intervals excluding zero through 18 months) and unemployment is lower (by 0.6 percentage points at twelve
months). After a rate \emph{decrease}, industrial production is lower by 1.4\% at twelve months and 2.2\% at
twenty-four, and unemployment is higher by 0.65 and 1.1 points. These are the opposite of the responses implied by
standard theory and found by most of the literature, and they are the same in the macro specification
(Appendix~\ref{app:robust}).

\paragraph{Staff forecasts remove the price puzzle.} The price responses behave differently. With the macro
information set, a rate increase is followed by significantly \emph{higher} prices from five months on, which is the
well-known price puzzle \citep{sims1992}. Conditioning on the Greenbook forecasts makes the price response to a rate increase
small and statistically indistinguishable from zero at every horizon (0.20\% at twelve months, interval
$[-0.07,0.39]$). This is the pattern \citet{romer2004} report: the puzzle reflects the Committee's reaction to
inflation it foresees, and the staff forecast captures that information. A rate decrease is followed by higher
prices under both specifications (0.8--1.0\% at twelve months), the direction theory predicts.

\begin{table}[!htbp]
  \centering
  \caption{Effects of rate changes relative to no change, Greenbook information set}\label{tab:results}
  \footnotesize
  \tablefile{results_main.tex}
  \par\medskip
  \begin{minipage}{0.86\linewidth}\footnotesize
  \emph{Notes:} Cumulative change $h$ months after the meeting. $\hat\theta_A$: effect on meetings whose menu contains
  both the action and no change, with basic moving-block bootstrap 90\% interval (499 draws, blocks of 8 meetings).
  Identified set: worst-case bounds on the average effect over all meetings (Proposition~\ref{prop:bounds}).
  $\lambda^\ast$: smallest heterogeneity between the non-overlap menus and the overlap region, in the outcome's
  units, that would make the sign of the average effect unidentified (Proposition~\ref{prop:lambda}).
  \end{minipage}
\end{table}

\paragraph{How much extrapolation is needed?} If effects were the same in every menu ($\lambda=0$), the
overlap-region estimates would be the average effects. The breakdown values in Table~\ref{tab:results} say how far
that homogeneity can be relaxed before the sign of the average effect is lost. For a rate increase, a third of the
meetings are outside the overlap region, so $\lambda^\ast=3\,|\hat\theta^{\cA}|$: the twelve-month IP response keeps its
sign unless the average effect in the menus without the status quo or without tightening differs from the
overlap-region effect by more than 5.7 percentage points. For a rate decrease two fifths of meetings lie outside the
region and $\lambda^\ast\approx2.4\,|\hat\theta^{\cA}|$. Figure~\ref{fig:sensitivity} traces the full set as a
function of $\lambda$. Reconciling the average effects with the textbook signs would require the effect in the
non-overlap menus to be not merely different from, but several times larger than and opposite in sign to, the effect
in the overlap region. That cannot be ruled out, but it would be an odd defence: the counter-theoretical signs come
from the meetings where the data are most informative, not from extrapolation.

\begin{figure}[!htbp]
  \centering
  \includegraphics[width=0.86\linewidth]{sensitivity_main.pdf}
  \caption{Identified set for the twelve-month effect as a function of $\lambda$, the maximal difference between
  the average effect in a non-overlap menu and in the overlap region (Assumption~\ref{as:lambda}). At $\lambda=0$ the set is
  the point $\hat\theta^{\cA}_d$; for large $\lambda$ it reaches the worst case. Greenbook propensity score.}\label{fig:sensitivity}
\end{figure}

\paragraph{Interpretation.} The overlap-region estimates are credible only under Assumption~\ref{as:unconf}.
Their signs indicate that it fails. The Committee tightens when it expects the economy to strengthen, and
the expected strength shows up in subsequent output whether or not the Committee acts. Neither lagged data, market
expectations nor the Greenbook forecasts absorb this completely. Two features of the sample make the problem
hard to escape. First, the identifying variation is thin: all but one of the 27 tightenings come from $\{U,T\}$
menus, so the effect of a rate increase is essentially a comparison of hikes and holds within those 43 meetings.
Second, weighting within a menu only reweights meetings the Committee itself judged similar enough to offer the
same options, so the remaining differences between hikes and holds are exactly the judgemental information that
is hardest to measure. Asymmetries between increases and decreases, which the 2020 version emphasised following
\citet{hamilton2002} and \citet{angrist2011}, cannot be separated from this selection problem with the present
data.
